\documentclass[11pt,a4paper]{report}

% --- Pacchetti Fondamentali ---
\usepackage[utf8]{inputenc}
\usepackage[T1]{fontenc}
\usepackage{lmodern}
\usepackage[italian]{babel}
\usepackage{etoolbox}
\usepackage[margin=1in]{geometry}
\usepackage{hyperref}
\usepackage{xcolor}
\usepackage{listings}
\usepackage{enumitem}
\usepackage{titlesec}
\usepackage{fancyhdr}
\usepackage{cleveref}
\usepackage{graphicx}
\usepackage{longtable}
\usepackage{booktabs}
\usepackage{setspace}

\crefformat{section}{\S#2#1#3}
\crefformat{subsection}{\S#2#1#3}
\crefformat{subsubsection}{\S#2#1#3}
\crefformat{figure}{#2Figura~#1#3}
\crefformat{footnote}{#2\footnotemark[#1]#3}
\crefformat{table}{#2Tabella~#1#3}

% --- Configurazione Header e Footer ---
\pagestyle{fancy}
\fancyhf{}
\fancyhead[L]{\small Programmazione e Tecnologie Web}
\fancyhead[R]{\small Giuseppe Capasso}
\fancyfoot[C]{\thepage}
\renewcommand{\headrulewidth}{0.4pt}
\usepackage{tikz}
\usetikzlibrary{arrows.meta, positioning, shapes.geometric, decorations.pathreplacing, calc, shadows, fit}

% --- Configurazione Colori e Link ---
\definecolor{primary}{RGB}{38, 66, 139}
\definecolor{codebg}{RGB}{245, 245, 245}
\definecolor{keywordcolor}{RGB}{0, 0, 255}
\definecolor{stringcolor}{RGB}{163, 21, 21}
\definecolor{commentcolor}{RGB}{0, 128, 0}

\hypersetup{
    colorlinks=true,
    linkcolor=primary,
    urlcolor=primary,
    citecolor=primary,
}

% --- Configurazione Formattazione Codice ---
\lstdefinelanguage{json}{
    basicstyle=\ttfamily\small,
    numbers=left,
    numberstyle=\scriptsize,
    stepnumber=1,
    numbersep=8pt,
    showstringspaces=false,
    breaklines=true,
    frame=lines,
    backgroundcolor=\color{codebg},
    string=[s]{"}{"},
    comment=[l]{//},
    morecomment=[s]{/*}{*/},
    stringstyle=\color{stringcolor},
    keywordstyle=\color{keywordcolor},
}

\lstdefinelanguage{javascript}{
    keywords={async, await, break, case, catch, class, const, continue, debugger, default, delete, do, else, export, extends, false, finally, for, function, if, import, in, instanceof, new, null, return, super, switch, this, throw, true, try, typeof, var, void, while, with, yield, let},
    morekeywords={console, log, Error, fetch, JSON, parse, stringify},
    sensitive=true,
    morecomment=[l]{//},
    morecomment=[s]{/*}{*/},
    morestring=[b]',
    morestring=[b]",
    morestring=[b]`,
    keywordstyle=\color{primary}\bfseries,
    commentstyle=\color{commentcolor}\itshape,
    stringstyle=\color{stringcolor},
    basicstyle=\ttfamily\small
}

\lstset{
    backgroundcolor=\color{codebg},
    basicstyle=\ttfamily\small,
    breaklines=true,
    frame=single,
    rulecolor=\color{lightgray},
    captionpos=b,
    xleftmargin=1em,
    xrightmargin=1em,
    keywordstyle=\color{keywordcolor},
    commentstyle=\color{commentcolor},
    stringstyle=\color{stringcolor}
}

% --- Formattazione Titoli ---
\titleformat{\chapter}[display]
  {\normalfont\bfseries\color{primary}}{\huge Capitolo \thechapter}{1ex}{\Huge}
\titlespacing*{\chapter}{0pt}{-20pt}{40pt}

\onehalfspacing

\begin{document}

% --- Frontespizio ---
\begin{titlepage}
    \centering
    \vspace*{4cm}
    {\Huge \textbf{\color{primary} Programmazione e Tecnologie Web}} \\[1.5cm]
    {\Large \textbf{Dispense del Corso - Lezione 4}} \\[0.5cm]
    {\large Docente: Giuseppe Capasso} \\[1.5cm]
    {\normalsize \textbf{Cloud Native Development \& Cybersecurity Specialist}} \\
    {\normalsize Istituto Tecnico Superiore (ITS)} \\
    \vfill
    {\normalsize Anno Accademico 2025/2026}
\end{titlepage}

\newpage
\tableofcontents
\newpage

% ==========================================
% CAPITOLO 1
% ==========================================
\chapter{Architetture Backend}

\section{Il Ruolo del Server nello Stack Web}
In un'architettura web moderna, il server agisce come nodo centrale per la gestione del ciclo di vita delle risorse. A differenza del frontend, che risiede nel contesto limitato e spesso insicuro del browser, il backend opera in un ambiente \textbf{Trusted} (fido). 

Il server-side ha accesso diretto alle risorse di calcolo, al file system e alle logiche di validazione dei dati che non possono essere delegate al client per motivi di sicurezza e integrità. La sua funzione primaria è quella di esporre interfacce (API) che permettano di interagire con le risorse applicative in modo sicuro, garantendo che le regole di business siano applicate uniformemente a tutti i client.

\subsection{L'Ambiente di Esecuzione (Runtime)}
Mentre il codice frontend viene eseguito nell'interprete JavaScript del browser (V8, SpiderMonkey), il backend vive in un sistema operativo (Linux, Windows, macOS). Questo comporta:
\begin{itemize}
    \item \textbf{Accesso alle Variabili d'Ambiente}: Gestione di segreti e configurazioni tramite file \texttt{.env}.
    \item \textbf{Persistenza}: Capacità di scrivere e leggere permanentemente su disco.
    \item \textbf{Networking}: Possibilità di contattare altri server, database o servizi esterni (es. gateway di pagamento).
\end{itemize}

\subsection*{Nota Metodologica: Persistenza su File e Focalizzazione Didattica}
Per le finalità di questo corso, che mira a fornire una solida comprensione delle dinamiche Web e dei protocolli di integrazione, abbiamo adottato una strategia di persistenza basata su file JSON. Questa scelta didattica evita l'introduzione precoce di sistemi complessi di gestione database (DBMS), permettendoci di dedicare l'intero carico didattico alla progettazione delle API, alla gestione del flusso HTTP e all'architettura dei sistemi distribuiti.

\section{Architettura degli Handler e Ciclo della Richiesta}
Il flusso di controllo nel backend è regolato dagli \textbf{Handler}. Un handler è una funzione che associa una determinata rotta (endpoint) a un metodo HTTP. Il server delega all'handler la logica necessaria per processare la richiesta, validare l'input ed eseguire l'operazione richiesta sulla risorsa.

\begin{figure}[ht]
\centering
\begin{tikzpicture}[
    request/.style={rectangle, draw, fill=blue!10, minimum width=2cm, minimum height=0.8cm},
    router/.style={diamond, draw, fill=yellow!10, minimum width=2.5cm, minimum height=1.5cm, align=center},
    handler/.style={rectangle, draw, fill=primary!10, minimum width=2cm, minimum height=0.8cm},
    file/.style={cylinder, draw, rotate=90, fill=gray!10, minimum width=0.8cm, minimum height=1cm}
]
    \node (req) [request] {HTTP Request};
    \node (rot) [router, right=1.5cm of req] {Router};
    \node (h1) [handler, right=1.5cm of rot, yshift=1cm] {Handler GET};
    \node (h2) [handler, right=1.5cm of rot, yshift=-1cm] {Handler POST};
    \node (f) [file, right=1.5cm of h2] {data.json};

    \draw [-latex] (req) -- (rot);
    \draw [-latex] (rot) -- node[above, scale=0.7] {/tasks} (h1);
    \draw [-latex] (rot) -- node[above, scale=0.7] {/tasks} (h2);
    \draw [-latex] (h2.east) -- (f.north);
\end{tikzpicture}
\caption{Diagramma di flusso: dalla richiesta al salvataggio su file via Handler.}
\end{figure}

Oltre al routing, l'handler gestisce internamente il ciclo di vita della richiesta:

\begin{figure}[ht]
\centering
\begin{tikzpicture}[
    block/.style={rectangle, draw, fill=blue!10, minimum width=3cm, minimum height=1cm, align=center},
    arrow/.style={-latex, thick}
]
    \node (val) [block] {1. Validazione\\Input};
    \node (exec) [block, right=0.8cm of val] {2. Esecuzione\\Logica};
    \node (resp) [block, right=0.8cm of exec] {3. Preparazione\\Risposta};

    \draw [arrow] (val) -- (exec);
    \draw [arrow] (exec) -- (resp);
\end{tikzpicture}
\caption{Ciclo di vita di una richiesta all'interno di un Handler.}
\end{figure}

\begin{itemize}[leftmargin=3em]
    \item \textbf{Validazione}: Verifica che l'input ricevuto dal client sia conforme al contratto dell'API (es. presenza campi obbligatori, tipi di dato corretti). Se la validazione fallisce, il ciclo si interrompe immediatamente restituendo un errore \texttt{400 Bad Request}.
    \item \textbf{Esecuzione}: Esecuzione delle operazioni logiche (es. modifica file, calcoli, interrogazione di altri servizi).
    \item \textbf{Risposta}: Costruzione del payload di risposta (solitamente in JSON) e definizione del codice di stato HTTP appropriato (\texttt{200}, \texttt{201}, ecc.).
\end{itemize}


% ==========================================
% CAPITOLO 2
% ==========================================
\chapter{Setup del Server}

\section{Gestione delle Dipendenze}
Nello sviluppo backend moderno, non si scrive tutto da zero. Ci si appoggia a framework e librerie gestite tramite \textbf{Package Manager}.
\begin{itemize}
    \item \textbf{Python}: Utilizza \texttt{pip} e i file \texttt{requirements.txt} (o \texttt{Poetry/Pyproject.toml}). \`E fondamentale l'uso di \textit{Virtual Environments} per isolare le dipendenze del progetto.
    \item \textbf{Node.js}: Utilizza \texttt{npm} o \texttt{yarn} e il file \texttt{package.json}. La cartella \texttt{node\_modules} contiene tutte le librerie installate.
\end{itemize}

\section{Python: Server con Flask}
Flask è un micro-framework Python estremamente popolare per la sua semplicità e la curva di apprendimento dolce. È ideale per prototipare rapidamente API REST.

\subsection{Installazione ed Esecuzione}
Assicuratevi di avere un ambiente virtuale attivo e installate Flask:
\begin{lstlisting}[language=bash]
pip install flask
\end{lstlisting}

Per avviare l'applicazione in modalità sviluppo, utilizzate il comando integrato. L'opzione \texttt{--debug} abilita l'\textbf{Hot Reloading}: il server si riavvia automaticamente ad ogni modifica del codice.
\begin{lstlisting}[language=bash]
export FLASK_APP=app.py
flask run --debug
\end{lstlisting}

\subsection{Codice Base}
Il seguente esempio inizializza un'applicazione Flask.

\begin{lstlisting}[language=python, caption=app.py (Flask)]
from flask import Flask, jsonify

app = Flask(__name__)

@app.route('/', methods=['GET'])
def index():
    return jsonify({"message": "Server Flask attivo!"})
\end{lstlisting}

\section{Node.js: Server con Express.js}
Express.js è lo standard de facto per le applicazioni web in Node.js. Offre un sistema di middleware flessibile che semplifica enormemente la gestione delle richieste.

\subsection{Setup del Progetto}
Inizializzate il progetto e installate Express:
\begin{lstlisting}[language=bash]
npm init -y
npm install express
\end{lstlisting}

\subsection{Sviluppo e Hot Reloading}
In ambito Node.js, si utilizza spesso \texttt{nodemon} per evitare di riavviare manualmente il server:
\begin{lstlisting}[language=bash]
npx nodemon index.js
\end{lstlisting}

\subsection{Codice Base}
Creiamo un server che esegue le stesse operazioni dell'esempio Flask.

\begin{lstlisting}[language=javascript, caption=index.js (Express)]
const express = require('express');
const app = express();
const port = 3000;

app.get('/', (req, res) => {
    res.json({ message: "Server Express attivo!" });
});

app.listen(port, () => {
    console.log(`Server in ascolto sulla porta ${port}`);
});
\end{lstlisting}

% ==========================================
% CAPITOLO 3
% ==========================================
\chapter{Implementazione API REST}

\section{Modellazione delle Risorse}
In un'API REST, tutto ruota attorno alle \textbf{Risorse}. Una risorsa è un'entità logica identificata da un URI (Uniform Resource Identifier). La modellazione corretta prevede che gli URI utilizzino sostantivi al plurale (es. \texttt{/tasks}) e che i verbi siano espressi unicamente dai metodi HTTP.

\section{Semantica dei Metodi e Codici di Stato}
\`E fondamentale che il server risponda con il codice di stato più appropriato per descrivere l'esito dell'operazione. Questo permette al client di gestire correttamente i diversi scenari (successo, errore di validazione, risorsa mancante).

\begin{table}[ht]
\centering
\begin{tabular}{lcl}
\toprule
\textbf{Operazione} & \textbf{Metodo} & \textbf{Codice di Stato Tipico} \\
\midrule
Lettura Lista       & \texttt{GET}    & \texttt{200 OK} \\
Creazione Risorsa   & \texttt{POST}   & \texttt{201 Created} \\
Aggiornamento       & \texttt{PUT}    & \texttt{200 OK} o \texttt{204 No Content} \\
Cancellazione       & \texttt{DELETE} & \texttt{200 OK} o \texttt{204 No Content} \\
Errore Validazione  & -               & \texttt{400 Bad Request} \\
Risorsa Inesistente & -               & \texttt{404 Not Found} \\
\bottomrule
\end{tabular}
\caption{Corrispondenza tra operazioni REST e codici di stato HTTP.}
\label{tab:rest_status_codes}
\end{table}

\section{Implementazione in Flask}
Flask utilizza il decoratore \texttt{@app.route} per mappare URL e metodi. Notate l'uso di \texttt{request.get\_json()} per estrarre il payload inviato dal client.

\begin{lstlisting}[language=python, caption=Implementazione REST in Flask]
from flask import Flask, request, jsonify

app = Flask(__name__)

tasks = []

@app.route('/tasks', methods=['GET'])
def get_tasks():
    return jsonify(tasks)

@app.route('/tasks', methods=['POST'])
def add_task():
    new_task = request.get_json()
    # Esempio di validazione minima
    if not new_task or 'title' not in new_task:
        return jsonify({"error": "Dati mancanti"}), 400
    
    tasks.append(new_task)
    return jsonify({"status": "created"}), 201
\end{lstlisting}

\section{Implementazione in Express.js}
In Express, i metodi sono direttamente accessibili tramite l'istanza \texttt{app}. \`E necessario abilitare il middleware \texttt{express.json()} per permettere al server di interpretare il corpo delle richieste in formato JSON.

\begin{lstlisting}[language=javascript, caption=Implementazione REST in Express]
const express = require('express');
const app = express();

// Middleware fondamentale per il parsing del body JSON
app.use(express.json());

let tasks = [];

app.get('/tasks', (req, res) => {
    res.json(tasks);
});

app.post('/tasks', (req, res) => {
    const newTask = req.body;
    
    if (!newTask.title) {
        return res.status(400).json({ error: "Titolo richiesto" });
    }
    
    tasks.push(newTask);
    res.status(201).json({ status: "created" });
});
\end{lstlisting}


% ==========================================
% CAPITOLO 4
% ==========================================
\chapter{Persistenza su File JSON}

\section{Strategie di Lettura e Scrittura}
Per mantenere lo stato delle risorse tra un riavvio del server e l'altro, persistiamo i dati in formato JSON. L'operazione di persistenza richiede una gestione cautelativa:
\begin{enumerate}
    \item \textbf{Lettura}: Caricamento del file in memoria all'avvio o alla ricezione di una richiesta.
    \item \textbf{Modifica}: Manipolazione del dato in memoria (es. \texttt{append} a un array).
    \item \textbf{Scrittura}: Serializzazione e salvataggio sul disco.
\end{enumerate}

\section{Sfide Tecniche: Atomicità e Concorrenza}
L'uso di file piatti (\textit{flat files}) per la persistenza introduce problemi critici in ambienti di produzione:
\begin{itemize}
    \item \textbf{Race Conditions}: Se due client inviano una \texttt{POST} simultaneamente, il server potrebbe leggere lo stesso stato del file, aggiungere entrambi i task in memoria, ma la seconda scrittura sovrascriverebbe la prima, causando la perdita di un dato.
    \item \textbf{Performance}: Ogni operazione richiede la ri-scrittura dell'intero file, il che diventa inefficiente con migliaia di record.
\end{itemize}

\subsection*{Verso il Database (DBMS)}
Per risolvere questi limiti, le applicazioni reali utilizzano i \textbf{Database} (come SQLite, PostgreSQL o MongoDB). I database garantiscono le proprietà \textbf{ACID} (Atomicità, Consistenza, Isolamento, Durabilità), permettendo accessi concorrenti sicuri e performance ottimizzate tramite indici.

\section{Persistenza in Flask}
Utilizziamo il modulo standard \texttt{json} di Python. \`E buona norma gestire le eccezioni (es. file non trovato).

\begin{lstlisting}[language=python, caption=Gestione JSON in Flask]
import json
import os

FILE_PATH = 'tasks.json'

def load_tasks():
    if not os.path.exists(FILE_PATH):
        return []
    with open(FILE_PATH, 'r') as f:
        return json.load(f)

def save_tasks(tasks):
    with open(FILE_PATH, 'w') as f:
        json.dump(tasks, f, indent=4)
\end{lstlisting}

\section{Persistenza in Express.js}
Utilizziamo il modulo \texttt{fs} di Node.js. In contesti didattici preferiamo le versioni sincrone (\texttt{readFileSync}), ma in produzione si utilizzano sempre le versioni asincrone per non bloccare il loop degli eventi.

\begin{lstlisting}[language=javascript, caption=Gestione JSON in Express]
const fs = require('fs');
const FILE_PATH = './tasks.json';

const loadTasks = () => {
    try {
        if (!fs.existsSync(FILE_PATH)) return [];
        const data = fs.readFileSync(FILE_PATH);
        return JSON.parse(data);
    } catch (err) {
        return [];
    }
};

const saveTasks = (tasks) => {
    fs.writeFileSync(FILE_PATH, JSON.stringify(tasks, null, 4));
};
\end{lstlisting}


% ==========================================
% CAPITOLO 5
% ==========================================
\chapter{Sicurezza Applicativa}

\section{Validazione e Sanitizzazione}
Nello sviluppo backend, l'input del client deve essere considerato intrinsecamente non sicuro. La protezione delle risorse si basa su due pilastri:
\begin{itemize}[leftmargin=3em]
    \item \textbf{Validazione}: Verifica della conformità del dato rispetto allo schema atteso (tipo di dato, lunghezza, presenza di campi obbligatori).
    \item \textbf{Sanitizzazione}: Rimozione o codifica di caratteri potenzialmente dannosi per prevenire attacchi di tipo Injection (es. XSS o SQL Injection).
\end{itemize}

\section{Il Middleware come Pipeline di Sicurezza}
Nelle architetture modulari, il \textbf{Middleware} funge da strato intermedio tra la ricezione della richiesta e l'esecuzione della logica di business. Questo approccio a "pipeline" permette di isolare le responsabilità di sicurezza, garantendo che ogni richiesta superi controlli sequenziali prima di accedere alle risorse.

\begin{figure}[ht]
\centering
\begin{tikzpicture}[
    req/.style={rectangle, draw, fill=blue!10, minimum width=2cm, minimum height=0.8cm},
    mid/.style={rectangle, draw, fill=red!10, minimum width=2.5cm, minimum height=0.8cm, align=center},
    handler/.style={rectangle, draw, fill=green!10, minimum width=2cm, minimum height=0.8cm},
    arrow/.style={-latex, thick}
]
    \node (r) [req] {Request};
    \node (m1) [mid, right=0.8cm of r] {Middleware:\\Autenticazione};
    \node (m2) [mid, right=0.8cm of m1] {Middleware:\\Validazione};
    \node (h) [handler, right=0.8cm of m2] {Handler};

    \draw [arrow] (r) -- (m1);
    \draw [arrow] (m1) -- (m2);
    \draw [arrow] (m2) -- (h);
\end{tikzpicture}
\caption{Architettura a Middleware: pipeline di filtraggio e validazione.}
\end{figure}

\section{CORS: Cross-Origin Resource Sharing}
Il \textbf{CORS} è un meccanismo di sicurezza implementato dai browser per limitare le richieste cross-origin. Senza un'esplicita autorizzazione del server tramite header HTTP specifici, un frontend servito da un'origine (es. \texttt{localhost:5500}) non può interagire con un backend situato su un'altra origine (es. \texttt{localhost:3000}).

Il server deve esporre l'header \texttt{Access-Control-Allow-Origin} per definire quali client sono autorizzati a consumare le API.

\section{Autenticazione HTTP Basic}
L'autenticazione Basic è un protocollo standard per la trasmissione di credenziali. Il client invia lo username e la password codificati in Base64 all'interno dell'header \texttt{Authorization}.

\subsection*{Implementazione in Flask}
\begin{lstlisting}[language=python, caption=Basic Auth in Flask]
from flask import request, Response

def check_auth(username, password):
    # In produzione, utilizzare variabili d'ambiente o database
    return username == 'admin' and password == 'secret'

@app.route('/secure')
def secure_endpoint():
    auth = request.authorization
    if not auth or not check_auth(auth.username, auth.password):
        return Response('Accesso negato', 401, 
                       {'WWW-Authenticate': 'Basic realm="Login"'})
    return "Contenuto protetto"
\end{lstlisting}

\subsection*{Implementazione in Express.js}
\begin{lstlisting}[language=javascript, caption=Basic Auth in Express]
const basicAuth = require('express-basic-auth');

app.use(basicAuth({
    users: { 'admin': 'secret' },
    challenge: true
}));

app.get('/secure', (req, res) => {
    res.send("Contenuto protetto");
});
\end{lstlisting}


% ==========================================
% CAPITOLO 6
% ==========================================
\chapter{Integrazione Frontend-Backend}

\section{Gestione Asincrona con Fetch API}
L'interazione con il backend avviene tramite la \texttt{fetch} API, che implementa il modello delle \textit{Promise}. L'utilizzo di \texttt{async/await} permette una gestione lineare e leggibile del ciclo di vita della richiesta.

\begin{lstlisting}[language=javascript, caption=Pattern di consumo API asincrono]
async function getTasks() {
    try {
        const response = await fetch('http://localhost:3000/tasks');
        
        if (!response.ok) throw new Error(`HTTP error! status: ${response.status}`);
        
        const data = await response.json();
        console.log('Data received:', data);
    } catch (error) {
        console.error('Request failed:', error);
    }
}
\end{lstlisting}

\section{UX e Gestione degli Stati della Richiesta}
Un'integrazione robusta deve comunicare lo stato dell'operazione all'utente finale per garantire una User Experience (UX) coerente.
\begin{itemize}[leftmargin=3em]
    \item \textbf{Loading State}: Gestione visiva dell'attesa durante la latenza di rete.
    \item \textbf{Error Handling}: Comunicazione chiara di eventuali fallimenti (es. \texttt{401 Unauthorized} o \texttt{500 Server Error}).
    \item \textbf{Success Feedback}: Conferma esplicita dell'avvenuto aggiornamento dello stato sul server.
\end{itemize}

\section{Autenticazione lato Client}
Il client deve gestire la codifica delle credenziali prima della trasmissione. La funzione \texttt{btoa()} è lo standard browser per la codifica Base64.

\begin{lstlisting}[language=javascript, caption=Fetch con Basic Auth dinamica]
const base64Credentials = btoa(`${username}:${password}`);

fetch('http://localhost:3000/secure', {
    method: 'GET',
    headers: {
        'Authorization': `Basic ${base64Credentials}`
    }
})
.then(response => {
    if (response.status === 401) throw new Error('Unauthorized');
    return response.text();
})
.then(data => console.log('Success:', data))
.catch(err => alert('Auth failed: ' + err.message));
\end{lstlisting}

\section{Conclusioni}
La progettazione di sistemi web moderni richiede una comprensione profonda del disaccoppiamento tra presentazione (Frontend) e logica di business (Backend). La robustezza di un'applicazione risiede nella capacità del backend di garantire l'integrità dei dati e nella capacità del frontend di gestire correttamente il flusso asincrono e la comunicazione degli stati. L'adozione di standard come REST e protocolli di sicurezza consolidati è fondamentale per la scalabilità e la manutenibilità del software.

\end{document}
